\documentclass{book}
\usepackage[utf8]{inputenc}
\usepackage[spanish]{babel}
\usepackage{color}
\usepackage{graphicx}
\usepackage{enumerate}
\usepackage{float}
\begin{document}
	\frontmatter
		\thispagestyle{empty}

\begin{minipage}{0.15\textwidth}
	\includegraphics[width=0.9\textwidth]{portada/ipn.eps}
%	\includegraphics[width=0.9\textwidth]{portada/ipn.jpg}
\end{minipage}%
\begin{minipage}{0.65\textwidth}
\begin{center}
\large \bf Instituto Politécnico Nacional \newline
\large \bf Escuela Superior de Cómputo \newline
Analisis y Diseño Orientado a Objetos
\end{center}
\end{minipage}
\begin{minipage}{0.20\textwidth}
	\includegraphics[width=1\textwidth]{portada/escom.eps}
%\includegraphics[width=1\textwidth]{portada/escom.jpg}

\end{minipage}%

\vspace{3cm}
\centerline{\Large \bf Urban Reyes Adan}


\vspace{2cm}
\centerline{\Large  \bf Resumen de Articulos}

\vspace{2cm}

\centerline{\Large Profra: M. en C. Tanibet Pérez de los Santos Mondragón }
\vspace{1.5cm}
\centerline{\Large  Primer Parcial }

 
\vspace{3cm}
{\large  M\'{e}xico, D.F. \hfill 2014}

		\begin{minipage}{0.15\textwidth}
	        	\includegraphics[width=7\textwidth]{cu.eps}
		\end{minipage}

	\begin{table}[H]
		\centering
		\begin{tabular}{p{3cm} p{11cm}}
			\hline
				Caso de Uso & Buscar medicamento\\
			\hline 
				Resumen & Este caso de uso descrive la forma en como el servidor busca un medicamento\\
			\hline
				PreCondicion&El usuario tuvo que aver ingresado al servidor web y solicitado alguna informacion\\
			\hline
				Actor & Servidor Web\\
			\hline
				Descripcion&\\
				Flujo Principal&el servidor prepara la consulta
					\\&el servidor ejecuta la consulta
					\\&el servidor recupera el resultado de la consulta
					\\&el servidor imprime el resultado de la consulta (codigo de barras, nombre,activo(s),docificacion,indicaciones medicas) del medicamento\\
			\hline
				Flujo Secundario&el servidor prepara la consulta\\
				"consulta no valida"&el servidor ejecuta la consulta
				\\&mysql lanza una excepcion de consulta no valida
				\\&el servidor imprime la excepcion\\
			\hline
				Flujo Secundario&el servidor prepara la consulta\\
				"ningun resultado"&el servidor ejecuta la consulta
				\\&mysql no regresa ningun resultado
				\\&el servidor imprime no hay resultados para la consulta\\
			\hline
				PosCondicion&el servidor regresa al menu principal\\
			\hline
		\end{tabular}
	\end{table}

	\begin{table}[H]
		\centering
		\begin{tabular}{p{3cm} p{11cm}}
			\hline
				Caso de Uso & Consultar Medicamento\\
			\hline 
				Resumen &  Mediante este caso de uso se podra consultar el codigo de barras de un medicamento asi como su nombre el o los activos que contiene, para que sirve el medicamento y la docificacion.\\
			\hline
				PreCondicion&que el Usuario aya ingresado ala pagina web de nuestro servidor\\
			\hline
				Actor & Cliente o Vendedor o Administrador\\
			\hline
				Descripcion&\\
 				Flujo principal&
					el usuario ingresa ala pagina web del servidor
					\\&el servidor le pide que proporcione el codigo de barras del medicamento o el nombre o las sales del medicamento
					\\&el usuario llena uno de los 3 campos
					\\&se ejecuta el caso de uso Buscar Medicamento\\
			\hline
				Flujo secundario & el usuario ingresa ala pagina web del servidor\\
				"pagina web incorrecta" & el navegador le indica que la pagina ingresada no existe\\
			\hline
				Flujo secundario & el usuario ingresa ala pagina web del servidor\\
				"el usuario no proporsiona ni un campo" & el servidor le pide que proporcione el codigo de barras del medicamento o el nombre o las sales del medicamento
					\\&el usuario no llena ni un campo
					\\&el servidor determina que no lleno ni un campo
					\\&el servidor Imprime un mensaje diciendo Deve llenar almenos un campo\\
			\hline
				Flujo secundario & el usuario ingresa ala pagina web del servidor\\
				"el medicamento no esta en la Base de Datos" & el servidor le pide que proporcione el codigo de barras del medicamento o el nombre o las sales del medicamento
					\\&el usuario llena uno de los 3 campos
					\\&se ejecuta el caso de uso Buscar Medicamento\\
			\hline
				PosCondicion&El usuario tubo que regresar al menu prinsipal de la pagina web\\
			\hline
		\end{tabular}
	\end{table}

	\begin{table}[H]
		\centering
		\begin{tabular}{p{3cm} p{11cm}}
			\hline
				Caso de Uso & Añadir medicamento ala lista\\
			\hline 
				Resumen & en este caso de uso se pueden añadir medicamentos ala base de datos\\
			\hline
				PreCondicion&haya una sesion activa\\
			\hline
				Actor & Vendedor o Administrador\\
			\hline
				Descripcion&\\
				Flujo Principal&el actor preciona el boton agregar medicamento
				\\&el servidor le pide que ingrese los campos codigo de barras,nombre,activo(s),distribuidor,precio de compra,precio de venta,existencia,docificacion,indicaciones medicas
				\\&el actor ingresa los campos solicitados y preciona el boton agregar
				\\&el servidor recupera los campos que el actor ingreso
				\\&el servidor guarda la informacion el la base de datos\\
			\hline
				Flujo Secundario&el actor preciona el boto agregar medicamento\\
				"error al recuperar los campos"&el servidor le pide que ingrese los campos codigo de barras,nombre,activo(s),distribuidor,precio de compra,precio de venta,existencia,docificacion,indicaciones medicas
				\\&el actor ingresa los campos solicitados y preciona el boton agregar
				\\&el servidor determina que los campos solicitados los a ingresado mal (en lugar de ingresar un entero ingreso una cadena)
				\\&el servidor imprime el nombre del campo y le dice que deve ser un numero entero\\
			\hline
				Flujo Secundario&el actor preciona el boto agregar medicamento\\
				"el actor no llena menos campos de los solicitados"&el servidor le pide que ingrese los campos codigo de barras,nombre,activo(s),distribuidor,precio de compra,precio de venta,existencia,docificacion,indicaciones medicas
				\\&el actor no ingresa todos los campos solicitados y preciona el boton agregar
				\\&el servidor determina que no a llenado todos los campos
				\\&el servidor imprime es nesesario llenar todos los campos\\
			\hline
				Flujo Secundario&el actor preciona el boton agregar medicamento\\
				"el medicamto ya existe en la base de datos"&el servidor le pide que ingrese los campos codigo de barras,nombre,activo(s),distribuidor,precio de compra,precio de venta,existencia,docificacion,indicaciones medicas
				\\&el actor ingresa los campos solicitados y preciona el boton agregar
				\\&el servidor recupera los campos que el actor ingreso
				\\&mysql lanza una excion llave duplicada
				\\&el servidor imprime la execion\\
			\hline
				PosCondicion&la base de datos se alla actualizado\\
			\hline
		\end{tabular}
	\end{table}

	\begin{table}[H]
		\centering
		\begin{tabular}{p{3cm} p{11cm}}
			\hline
				Caso de Uso & Eliminar medicamento de la lista\\
			\hline 
				Resumen & En este caso de uso se eliminara un medicamento de la base de datos\\
			\hline
				PreCondicion&haya una sesion activa y el actor aya precionado elimar medicameto\\
			\hline
				Actor & Vendedor o Administrador\\
			\hline
				Descripcion&\\
				Flujo Principal&el servidor le pide al actor los siguientes campos codigo de barras,distribuidor del medicamento
				\\&el actor ingresa los campos solicitados
				\\&el servidor recupera los campos
				\\&el servidor elimina el medicamento de la base de datos\\
			\hline
				Flujo Secundario&el servidor le pide que ingrese los campos codigo de barras,distribuidor del medicamento\\
				"error al recuperar los campos"& el actor ingresa los campos solicitados y preciona el boton eliminar
				\\&el servidor determina que los campos solicitados los a ingresado mal (en lugar de ingresar un entero ingreso una cadena)
				\\&el servidor imprime el nombre del campo y le dice que deve ser un numero entero\\
			\hline
				Flujo Secundario&el servidor le pide que ingrese los campos codigo de barras, distribuidor del medicamento\\
				"el actor no llena menos campos de los solicitados"&el actor no ingresa todos los campos solicitados y preciona el boton eliminar
				\\&el servidor determina que no a llenado todos los campos
				\\&el servidor imprime es nesesario llenar todos los campos\\
			\hline
				Flujo Secundario&el servidor le pide que ingrese los campos codigo de barras, distribuidor del medicamento\\
				"el medicamto no existe en la base de datos"&el actor ingresa los campos solicitados y preciona el boton eliminar
				\\&el servidor recupera los campos que el actor ingreso
				\\&mysql lanza una excion no existe el registro
				\\&el servidor imprime la execion\\
			\hline
				PosCondicion&la base de datos tiene que estar actualizada\\
			\hline
		\end{tabular}
	\end{table}

	\begin{table}[H]
		\centering
		\begin{tabular}{p{3cm} p{11cm}}
			\hline
				Caso de Uso & Modificar plantilla\\
			\hline 
				Resumen & se modifican los campos del medicamento\\
			\hline
				PreCondicion& aya una sesion activa del actor y se alla precionado el boton modificar plantilla\\
			\hline
				Actor & Administrador\\
			\hline
				Descripcion&\\
				Flujo Principal&el servidor le pide que ingrese el nombre del campo a eliminar o añadir
					\\&el actor ingresa el nombre de un campo y presiona un boton
					\\&el servidor recupera el nombre del campo
					\\&el servidor ejecuta la consulta de acuerdo con el boton pulsado\\
			\hline
				Flujo Secundario&el servidor le pide que ingrese el nombre del campo a elimnar o añadir\\
					"no lleno el campo"&el actor no ingresa el campo y presiono uno de los botones
					\\&el servidor imprime el mensaje es nesasario ingresara el campo\\
			\hline
				Flujo Secundario&el servidor le pide que ingrese el nombre del campo a eliminar o añadir\\
					"el campo ya existe"&el actor ingresa el campo y presiona añadir	
					\\&el servidor recupera el campo ejecuta la consulta
					\\&mysql regresa una exepcion el campo existe ya
					\\&el servidor imprime la excepcion\\
			\hline
				Flujo Secundario&el servidor le pide que ingrese el nombre del campo a eliminar o añadir\\
					"el campo es parte de la llave"&el actor ingresa el campo y presiona eliminar
					\\&el servidor recupera el campo y ejecuta la consulta
					\\&mysql regresa una excepcion el campo es una llave
					\\&el servidor imprime la excepcion\\
			\hline
				PosCondicion&base de datos actualizada\\
			\hline
		\end{tabular}
	\end{table}

	\begin{table}[H]
		\centering
		\begin{tabular}{p{3cm} p{11cm}}
			\hline
				Caso de Uso & Añadir usuario\\
			\hline 
				Resumen & se añade un vendedor ala vase de datos\\
			\hline
				PreCondicion&Sesion activa del Administrador aver presionado gestionar Vendedores y despues presionado añadir\\
			\hline
				Actor & Administrador\\
			\hline
				Descripcion&\\
				Flujo Principal&el servidor le pide el nombre, contraseña y farmacia del nuevo vendedor
					\\&el actor ingresa los campos solicitados y presiona listo
					\\&el servidor ejecuta la consulta ala vase de datos\\
			\hline
				Flujo Secundario&el servidor le pide el nombre, contraseña y farmacia del nuevo vendedor\\
				"el vendedor ya existe"&el actor ingresa los campos solicitados y presiona listo
				\\&el servidor ejecuta la consulta
				\\&mysql lansa una excepcion el registro ya existe
				\\&el servidor imprime la excepcion\\
			\hline
				Flujo Secundario&el servidor le pide el nombre, contraseña y farmacia del nuevo vendedor\\
				"no lleno todos los campos"&el actor no ingresa todos los campos solicitados y presiona listo
				\\&el servidor determina que el actor no ingreso todos los campos
				\\&el servidor imprime deve llenar todos los campos\\
			\hline
				PosCondicion&averse añadido un registro ala tabla vendedores\\
			\hline
		\end{tabular}
	\end{table}

	\begin{table}[H]
		\centering
		\begin{tabular}{p{3cm} p{11cm}}
			\hline
				Caso de Uso & Eliminar usuario\\
			\hline 
				Resumen & se elimina un vendedor ala vase de datos\\
			\hline
				PreCondicion&Sesion activa del Administrador aver presionado gestionar Vendedores y despues presionado eliminar\\
			\hline
				Actor & Administrador\\
			\hline
				Descripcion&\\
				Flujo Principal&el servidor le pide el nombre y farmacia del vendedor
					\\&el actor ingresa los campos solicitados y presiona listo
					\\&el servidor ejecuta la consulta ala vase de datos\\
			\hline
				Flujo Secundario&el servidor le pide el nombre y farmacia del vendedor\\
				"el vendedor no existe"&el actor ingresa los campos solicitados y presiona listo
				\\&el servidor ejecuta la consulta
				\\&mysql lansa una excepcion el registro no existe
				\\&el servidor imprime la excepcion\\
			\hline
				Flujo Secundario&el servidor le pide el nombre y farmacia del vendedor\\
				"no lleno todos los campos"&el actor no ingresa todos los campos solicitados y presiona listo
				\\&el servidor determina que el actor no ingreso todos los campos
				\\&el servidor imprime deve llenar todos los campos\\
			\hline
				PosCondicion&averse eleminado un registro de la tabla vendedores\\
			\hline
		\end{tabular}
	\end{table}
	
	\begin{table}[H]
		\centering
		\begin{tabular}{p{3cm} p{11cm}}
			\hline
				Caso de Uso & Modificar usuario\\
			\hline 
				Resumen & se modifica el nombre o contraseña o farmacia de un vendedor en la base de datos\\
			\hline
				PreCondicion&Sesion activa del Administrador aver presionado gestionar Vendedores y despues presionado modificar\\
			\hline
				Actor & Administrador\\
			\hline
				Descripcion&\\
				Flujo Principal&el servidor le pide el nombre y farmacia del vendedor
					\\&el actor ingresa los campos solicitados y presiona listo
					\\&el servidor ejecuta la consulta ala vase de datos
					\\&el servidor muestra la informacio y le pide que modifique el campo que desee
					\\&el actor modifica los campos y presional listo 
					\\&el servidor guarda la informacion\\
			\hline
				Flujo Secundario&el servidor le pide el nombre y farmacia del vendedor\\
				"el vendedor ya existe"&el actor ingresa los campos solicitados y presiona listo
				\\&el servidor ejecuta la consulta
				\\&el servidor muestra la informaicon y le pide que modifique el campo que desee
				\\&el actor modifica los campos y presiona listo
				\\&el servidor guarda la informacion
				\\&mysql determina que el vendedor esta repetido y lanza una excepcion 
				\\&el servidor imprime la excepcion\\
			\hline
				Flujo Secundario&el servidor le pide el nombre y farmacia del vendedor\\
				"no lleno todos los campos"&el actor no ingresa todos los campos solicitados y presiona listo
				\\&el servidor determina que el actor no ingreso todos los campos
				\\&el servidor imprime deve llenar todos los campos\\
			\hline
				PosCondicion&se modifico registro ala tabla vendedores\\
			\hline
		\end{tabular}
	\end{table}

\end{document}
